\section{Why 4\% Network Reduction Matters at Scale}
At 68{,}000 cross-node operations per second, a 4.1\% reduction is $\sim$2{,}800 operations per second eliminated --- equivalent to roughly one Redis shard's worth of capacity, freeing it for growth without provisioning new infrastructure. The latency improvements compound: lower Redis load reduces Redis tail latency, which feeds back as further API latency reduction (the ``virtuous cycle'' visible in the 26.2\% Redis p99 drop).

\section{Trade-offs}

\subsection{Memory vs.\ Latency}
The two-tier cache trades $<$20\,MB per node for tens of microseconds saved per cached lookup. At 3{,}000 RPS per node this saves cumulative seconds of CPU per minute.

\subsection{Staleness vs.\ Freshness}
TTL-based caching admits bounded staleness; the cache-safety boundary framework formalises which data classes can tolerate it. The contention layer ensures that single-resource state --- which cannot tolerate staleness --- is never served from a stale cache.

\subsection{Strong vs.\ Eventual Consistency}
We use strong consistency (Lua, lock, OCC) for inventory and balances, and eventual consistency (PN-Counter) only for soft holds and analytics counters. The decision matrix (Table~\ref{tab:matrix}) operationalises this choice.

\section{Limitations}
\begin{itemize}[leftmargin=*]
    \item Distributed locks (Technique~1) remain vulnerable to misconfigured TTLs and clock drift; we mitigate with fence tokens but do not eliminate.
    \item PN-Counter over-allocation, while bounded, is non-zero and must be acceptable to the business owner of the workload.
    \item Pub/Sub-based invalidation has a tail latency window of tens of milliseconds during which the server-local cache may serve a freshly-stale display value. This is acceptable for display reads; authoritative reads bypass the cache.
\end{itemize}

% =============================================================================
